\chapter{分子对称性与正规振动}
\label{ch:molecular-symmetry}

正规振动最终来自质量加权 Hessian 的本征矢量，但在真正对角化一个 \(3N\times3N\) 矩阵之前，对称性已经能够回答大量问题：有多少种不同频率、哪些模式简并、哪些模式能够被红外或 Raman 光谱看见，以及 Hessian 可以分成哪些彼此不耦合的块。这些结论来自有限群表示的特征标理论，不必从若干分子的已知模式表反推一般规律。
\section{有限群作用、投影算符与对称分块}

\begin{definition}[有限群、表示与特征标]
有限群 \(G\) 是有限多个可复合、可逆的操作组成的集合。单位元记为 \(e\)，每个 \(g\) 的逆元记为 \(g^{-1}\)。群在复向量空间 \(V\) 上的表示，是满足
\[
D(e)=I,\qquad D(g_1g_2)=D(g_1)D(g_2)
\]
的可逆线性算符族。函数
\(\chi(g)=\operatorname{Tr}D(g)\) 称为这个表示的特征标。
\end{definition}

若 \(h=k g k^{-1}\)，则 \(g,h\) 共轭。表示矩阵满足
\(D(h)=D(k)D(g)D(k)^{-1}\)，相似矩阵的迹相等，所以
\(\chi(h)=\chi(g)\)。特征标在每个共轭类上为常数，计算群平均时才可以把逐元素求和改成逐共轭类求和。

\begin{definition}[不可约表示]
若子空间 \(W\subseteq V\) 满足 \(D(g)W\subseteq W\) 对所有
\(g\in G\) 成立，就称 \(W\) 为不变子空间。只有 \(0\) 与 \(V\) 两个不变子空间的非零表示称为不可约表示。
\end{definition}

在证明一般定理以前，先让这些定义在一个真实分子上工作。取水分子的平衡构型，把氧原子放在原点，分子平面取为 \(yz\) 平面，\(z\) 轴沿键角平分线。保持这一构型不变的操作有
\[
E,\qquad C_2(z),\qquad \sigma_{xz},\qquad \sigma_{yz}.
\]
其中 \(C_2(z)\) 把空间绕 \(z\) 轴转过 \(\pi\)，
\(\sigma_{xz}\) 与 \(\sigma_{yz}\) 分别关于两个平面反射。每个非单位操作做两次都回到 \(E\)，而任意两个不同的非单位操作相乘得到第三个。例如
\(C_2\sigma_{xz}=\sigma_{yz}\)。这四个操作组成点群 \(C_{2v}\)。

这些操作先作用在普通坐标向量
\((x,y,z)^{\mathsf T}\) 上。按上述坐标约定，表示矩阵为
\[
\begin{aligned}
R(E)&=\operatorname{diag}(1,1,1),\\
R(C_2)&=\operatorname{diag}(-1,-1,1),\\
R(\sigma_{xz})&=\operatorname{diag}(1,-1,1),\\
R(\sigma_{yz})&=\operatorname{diag}(-1,1,1).
\end{aligned}
\]
直接相乘便能核对
\(R(g_1g_2)=R(g_1)R(g_2)\)。这一步展示了“群”和“表示”的差别：群元是几何操作，矩阵是这些操作在所选向量空间和基底中的写法。若改用分子轨道或原子位移作为基，群没有改变，表示矩阵却会改变维数和具体形式。

由于 \(C_{2v}\) 是交换群，它的复不可约表示都是一维的。一维表示把每个群元送到数 \(\pm1\)，并且仍要保持乘法。以四个群元按
\((E,C_2,\sigma_{xz},\sigma_{yz})\) 排列，四个可能的特征标行是
\begin{equation}
\begin{array}{c|rrrr|l}
 &E&C_2&\sigma_{xz}&\sigma_{yz}&\text{常用基函数}\\ \hline
A_1&1& 1& 1& 1&z,\ x^2,y^2,z^2\\
A_2&1& 1&-1&-1&R_z,\ xy\\
B_1&1&-1& 1&-1&x,\ R_y,\ xz\\
B_2&1&-1&-1& 1&y,\ R_x,\ yz
\end{array}
\label{eq:water-c2v-character-table}
\end{equation}
例如 \(y\) 在 \(C_2\) 与 \(\sigma_{xz}\) 下变号，在
\(\sigma_{yz}\) 下不变，所以它属于 \(B_2\)。表中转动
\(\mathbf R\) 是轴矢量；反射下它比普通坐标多乘一个行列式
\(-1\)，因此不能照抄 \(x,y,z\) 的行。

现在计算水分子的九维微位移表示。一个群元若把某原子换到另一原子，该原子对应的三维位移块落在总矩阵的非对角位置，对迹没有贡献。只有原地不动的原子才贡献相应三维矩阵的迹。

单位元固定三个原子，每个贡献 \(3\)，所以
\(\chi_{3N}(E)=9\)。\(C_2\) 交换两个氢，只固定氧，而
\(\Tr R(C_2)=-1\)，故
\(\chi_{3N}(C_2)=-1\)。反射
\(\sigma_{xz}\) 也交换两个氢，只固定氧，三维反射矩阵的迹为
\(1\)，故特征标为 \(1\)。整个分子位于 \(yz\) 平面内，所以
\(\sigma_{yz}\) 固定全部三个原子，每个贡献 \(1\)，故特征标为
\(3\)。因此
\begin{equation}
\chi_{3N}=(9,-1,1,3).
\label{eq:water-3n-character}
\end{equation}

九个自由度中还包含整体平移和整体转动。平移与普通向量
\((x,y,z)\) 同变换，所以
\[
\chi_{\rm trans}=(3,-1,1,1).
\]
转动按轴矢量变换。真旋转下与普通向量相同，反射下多一个
\(\det R=-1\)，因此
\[
\chi_{\rm rot}=(3,-1,-1,-1).
\]
从九维位移空间中去掉这六个刚体方向，振动表示的特征标为
\begin{equation}
\chi_{\rm vib}
=\chi_{3N}-\chi_{\rm trans}-\chi_{\rm rot}
=(3,1,1,3).
\label{eq:water-vibrational-character}
\end{equation}

对四元素群，某个不可约表示出现次数就是对应两行逐项相乘后除以四。于是
\[
\begin{aligned}
n_{A_1}&=\tfrac14(3+1+1+3)=2,\\
n_{A_2}&=\tfrac14(3+1-1-3)=0,\\
n_{B_1}&=\tfrac14(3-1+1-3)=0,\\
n_{B_2}&=\tfrac14(3-1-1+3)=1.
\end{aligned}
\]
水分子的三个振动模式因此分解为
\begin{equation}
\Gamma_{\rm vib}=2A_1\oplus B_2.
\label{eq:water-mode-decomposition}
\end{equation}
两个 \(A_1\) 模是对称伸缩与弯曲，\(B_2\) 模是不对称伸缩。这个计算没有先求 Hessian 的本征矢，却已经得到模式数目、简并度和对称性。后面的一般特征标正交定理，正是把这里“逐行相乘再平均”的算法证明为任意有限群都成立。

\begin{theorem}[有限群特征标正交关系]
设 \(\chi_\alpha,\chi_\beta\) 是两个不等价复不可约表示的特征标，则
\[
\frac1{|G|}\sum_{g\in G}\chi_\alpha(g)^*\chi_\beta(g)
=\delta_{\alpha\beta}.
\]
\end{theorem}

\begin{proof}
先把有限群表示酉化：从任意内积出发定义
\[
\langle u,v\rangle_G
=\frac1{|G|}\sum_{g\in G}\langle D(g)u,D(g)v\rangle.
\]
变量替换 \(g\mapsto gh\) 表明这个内积在群作用下不变，因此可在正交归一基中把所有 \(D(g)\) 写成酉矩阵。

对任意线性映射 \(A:V_\alpha\to V_\beta\)，作群平均
\[
\overline A=\frac1{|G|}\sum_{g\in G}
D_\beta(g)A D_\alpha(g)^{-1}.
\]
换元可验证
\(D_\beta(h)\overline A=\overline A D_\alpha(h)\)，所以
\(\overline A\) 是交织算符。Schur 引理说明：当
\(\alpha\ne\beta\) 时 \(\overline A=0\)；当
\(\alpha=\beta\) 时 \(\overline A=(\operatorname{Tr}A/d_\alpha)I\)。
依次令 \(A\) 为矩阵单位 \(E_{ij}\)，取相应矩阵元并对对角指标求和，便得到上式。整个证明只用群平均、Schur 引理和迹的线性性。
\end{proof}

因此，若可约表示的特征标为 \(\chi\)，并有
\(\chi=\sum_\alpha n_\alpha\chi_\alpha\)，两边与
\(\chi_\beta\) 作群平均内积，立刻得到
\[
n_\beta=\frac1{|G|}\sum_{g\in G}\chi_\beta(g)^*\chi(g).
\]
这就是稍后使用的约化公式；\(n_\beta\) 必为非负整数，因为它数的是不可约子空间的拷贝数。

一个含 \(N\) 个原子的分子，在平衡构型附近有 \(3N\) 个笛卡尔微位移自由度。把第 \(a\) 个原子的位移写成

\[
\delta\mathbf r_a\in E_a\cong\mathbb R^3,
\]

则总位移空间是直和

\begin{equation}
\mathcal V_{3N}
=\bigoplus_{a=1}^{N}E_a.
\label{eq:molecular-displacement-representation}
\end{equation}

设分子的点群为有限群 \(G\)。任一 \(g\in G\) 同时做两件事：它把某个原子位置 \(a\) 置换到 \(g a\)，并用三维正交矩阵 \(R(g)\) 旋转或反射该原子的位移向量。因此

\[
D_{3N}(g):\mathcal V_{3N}\to\mathcal V_{3N}
\]

是一个线性表示，而且群乘法满足

\[
D_{3N}(g_1g_2)=D_{3N}(g_1)D_{3N}(g_2).
\]

所以分子振动中的“群论”不是额外贴在 Hessian 上的分类标签，而是平衡构型的几何对称性本来就在位移 Hilbert 空间上诱导了一个表示。

把 \(D_{3N}(g)\) 看成按原子分块的 \(N\times N\) 块矩阵，每个块是 \(3\times3\)。若 \(g\) 把原子 \(a\) 移到另一个原子 \(b\neq a\)，则相应的 \(3\times3\) 线性映射位于非对角块 \((b,a)\)，对总矩阵迹没有贡献。只有满足 \(ga=a\) 的不动原子会留下对角块 \(R_a(g)\)。因此

\begin{equation}
\chi_{3N}(g)
=\operatorname{Tr}D_{3N}(g)
=\sum_{a\in\mathrm{Fix}(g)}
\operatorname{Tr}R_a(g).
\label{eq:fixed-atom-character}
\end{equation}

这就是课堂笔记中“只数不动原子”的严格来源。对三维真旋转 \(C_n^m\)，转角 \(\theta=2\pi m/n\)，其三个本征值为 \(1,\ee^{\ii\theta},\ee^{-\ii\theta}\)，所以每个不动原子贡献

\[
\operatorname{Tr}R(C_n^m)
=1+2\cos\theta.
\]

对镜面反射，三个本征值是 \(1,1,-1\)，因此每个不动原子贡献 \(1\)。反演的三个本征值全为 \(-1\)，每个不动原子贡献 \(-3\)。这些常用“计数公式”都只是 \eqnrefs{eq:fixed-atom-character} 的特例。

课堂中也可以把 \(\{E_a\}\) 看成有限离散 base 上的向量丛。这个几何图像是合法的，但在普通正规振动问题里真正参与计算的是截面空间

\[
\Gamma(E)\cong\bigoplus_a E_a=\mathcal V_{3N}
\]

上的有限维群表示，而不需要连续流形上的联络、曲率等额外结构。

设 \(\Gamma\) 是 \(G\) 的一个可约表示，特征标为 \(\chi(g)\)。若它分解为

\[
\Gamma\cong\bigoplus_\alpha n_\alpha\Gamma^{(\alpha)},
\]

则

\[
\chi(g)=\sum_\alpha n_\alpha\chi_\alpha(g).
\]

利用有限群不可约特征标的正交关系

\[
\frac1{|G|}
\sum_{g\in G}
\chi_\alpha^*(g)\chi_\beta(g)
=\delta_{\alpha\beta},
\]

两边乘 \(\chi_\beta^*(g)\) 后对群求和，得到不可约表示 \(\Gamma^{(\beta)}\) 的重数

\begin{equation}
n_\beta
=\frac1{|G|}
\sum_{g\in G}
\chi_\beta^*(g)\chi(g)
=\frac1{|G|}
\sum_{C}|C|\,\chi_\beta^*(C)\chi(C).
\label{eq:character-reduction-formula}
\end{equation}

第二个等号把求和改写为共轭类之和。这个公式是后面所有“把 \(3N\) 位移表示分解成正常模式”的计算核心。

更强的对象是不可约投影算符。若 $d_\alpha$ 是不可约表示 $\alpha$ 的维数，则
\begin{equation}
P_\alpha
=\frac{d_\alpha}{|G|}
\sum_{g\in G}\chi_\alpha^*(g)D_{3N}(g).
\label{eq:molecular-irrep-projector}
\end{equation}
它把任意笛卡尔位移直接投影到 $\alpha$-等型子空间。若质量加权 Hessian $F$ 与平衡构型的全部对称操作相容，则
\begin{equation}
[F,D_{3N}(g)]=0\quad(\forall g\in G)
\quad\Longrightarrow\quad
[F,P_\alpha]=0.
\label{eq:hessian-symmetry-commutator}
\end{equation}
因此 Hessian 在不同不可约表示之间没有矩阵元，正规模问题可以先按对称性分块再数值对角化。

\begin{proof}[解]
由\eqnrefs{eq:molecular-irrep-projector}，
\begin{equation*}
 P_\alpha
 =\frac{d_\alpha}{|G|}
 \sum_{g\in G}\chi_\alpha^*(g)D(g),
\end{equation*}
若 $[F,D(g)]=0$ 对每个 $g\in G$ 成立，由换位子的线性性，
\begin{align*}
 [F,P_\alpha]
 &=F P_\alpha-P_\alpha F\\
 &=\frac{d_\alpha}{|G|}\sum_{g\in G}\chi_\alpha^*(g)\!\left(FD(g)-D(g)F\right)\\
 &=\frac{d_\alpha}{|G|}\sum_{g\in G}\chi_\alpha^*(g)[F,D(g)]=0.
\end{align*}
另一方面由有限群特征标正交关系
\begin{equation*}
 \frac1{|G|}\sum_{g\in G}\chi_\alpha^*(g)\chi_\beta(g)
 =\delta_{\alpha\beta},
\end{equation*}
以及表示的直和分解 $D(g)=\bigoplus_\beta n_\beta D_\beta(g)$，得到投影代数
\begin{equation*}
 P_\alpha P_\beta=\delta_{\alpha\beta}P_\alpha,
 \qquad
 \sum_\alpha P_\alpha=I,
\end{equation*}
即 $\{P_\alpha\}$ 是一组正交完备投影（这里 $P_\alpha$ 投影到整个等型分量；若同一 irrep 出现多重拷贝，还需在该块内继续对角化）。由完备性，
\begin{equation*}
 F=IF I
 =\!\sum_{\alpha,\beta}P_\alpha F P_\beta.
\end{equation*}
又由 $[F,P_\alpha]=0$ 即 $F P_\alpha=P_\alpha F$，对任意 $\alpha,\beta$，
\begin{equation*}
 P_\alpha F P_\beta
 =P_\alpha P_\beta F
 =\delta_{\alpha\beta}P_\alpha F,
\end{equation*}
交叉块 $P_\alpha F P_\beta$ 在 $\alpha\neq\beta$ 时为零，于是
\begin{equation*}
 F=\sum_\alpha P_\alpha F P_\alpha.
\end{equation*}
故"不同对称性的正常模不混合"不是经验规则，而是 Schur 分块在 Hessian 上的直接结果。

如水分子 $C_{2v}$ 群的 3 个振动模分属 $A_1$ 与 $B_2$，Hessian 自动分块；$\mathrm{CH_4}$ 的 $T_2$ 三重简并模由 Schur 分块保证不分裂。
\end{proof}

\section{Hessian 的对称分块与正规振动}
对非线性分子，整体平移给出 3 个自由度，整体转动再给出 3 个自由度。于是

\begin{equation}
\Gamma_{\rm vib}
=\Gamma_{3N}
-\Gamma_{\rm trans}
-\Gamma_{\rm rot},
\label{eq:vibrational-representation-subtraction}
\end{equation}

总维数为 \(3N-6\)。线性分子绕分子轴的“转动”不改变核构型，因此只有两个独立整体转动，自由度为 \(3N-5\)。

这里的减法是表示环中的形式差；实际含义是从 \(\mathcal V_{3N}\) 中取出与整体质心平移和刚体转动对应的不变子空间，剩余正交补空间才由内部振动占据。

更精确地，正规模问题应在质量加权坐标中完成。令平衡位置相对于质心为 \(\mathbf R_a\)，质量为 \(m_a\)，并定义
\[
\mathbf q_a=\sqrt{m_a}\,\delta\mathbf r_a.
\]
三个整体平移方向可由
\[
(\mathbf t_\mu)_a=\sqrt{m_a}\,\mathbf e_\mu,
\qquad \mu=x,y,z,
\]
张成；三个刚体转动方向可由
\[
(\mathbf r_\mu)_a=\sqrt{m_a}\,(\mathbf e_\mu\times\mathbf R_a)
\]
张成。对非线性分子，这六个向量在去除线性相关后给出平移--转动子空间 \(\mathcal V_{\rm ext}\)，其质量度量正交补
\[
\mathcal V_{\rm vib}=\mathcal V_{3N}\ominus\mathcal V_{\rm ext}
\]
才是内部振动空间。在线性分子中，绕分子轴的 \(\mathbf r_\mu\) 恒为零，因此只需扣除两个转动方向。

\begin{proof}[解]
设势能只依赖内部相对构型，平衡点为 \(\{\mathbf R_a\}\)。整体平移 \(\mathbf R_a\mapsto\mathbf R_a+\boldsymbol\epsilon\) 不改变势能，所以
\[
U(\{\mathbf R_a+\boldsymbol\epsilon\})=U(\{\mathbf R_a\}).
\]
对 \(\epsilon_\mu\) 求一次导数得到
\[
\sum_a\frac{\partial U}{\partial r_{a\mu}}=0.
\]
在平衡点再对 \(r_{b\nu}\) 求导，若记笛卡尔 Hessian 为
\[
K_{b\nu,a\mu}
=\frac{\partial^2U}{\partial r_{b\nu}\partial r_{a\mu}},
\]
则
\[
\sum_a K_{b\nu,a\mu}=0.
\]
质量加权 Hessian 定义为
\[
F_{b\nu,a\mu}
=\frac{K_{b\nu,a\mu}}{\sqrt{m_bm_a}}.
\]
因此
\[
(F\mathbf t_\mu)_{b\nu}
=\sum_a\frac{K_{b\nu,a\mu}}{\sqrt{m_bm_a}}\sqrt{m_a}
=\frac1{\sqrt{m_b}}\sum_a K_{b\nu,a\mu}=0,
\]
即 \(F\mathbf t_\mu=0\)。

对无外场的孤立分子，整体转动同样不改变势能。令无穷小转角为 \(\delta\boldsymbol\theta\)，则
\[
\delta\mathbf R_a
=\delta\boldsymbol\theta\times\mathbf R_a.
\]
对旋转不变性做同样的二阶微分，在平衡点得到
\[
F\mathbf r_\mu=0.
\]
因此平移和刚体转动不是“经验上需要减掉的自由度”，而是由 Euclidean 对称性产生的 Hessian 零方向。真正的振动频率来自 \(F\) 在 \(\mathcal V_{\rm vib}\) 上的正谱。

实例。非线性分子有 3 个平移零模 + 3 个转动零模 = 6 个零频，$3N-6$ 个振动模。线性分子只有 2 个转动零模（绕分子轴转动无惯量），故 $3N-5$ 个振动模。水分子 $N=3$：$9-6=3$ 振动模。二氧化碳 $N=4$（线性）：$12-5=7$ 振动模。零模的对称性分类在群论约化中自动处理，确保振动表示 $\Gamma_{\rm vib}=\Gamma_{\rm tot}-\Gamma_{\rm trans}-\Gamma_{\rm rot}$ 的正确性。
\end{proof}

CH\(_4\) 的平衡构型属于 \(T_d\) 点群，群阶为 \(|T_d|=24\)，共轭类按

\[
E,
\quad 8C_3,
\quad 3C_2,
\quad 6S_4,
\quad 6\sigma_d
\]

排列。分子有一个中心 C 和四个 H，因此 \(3N=15\)。用 \eqnrefs{eq:fixed-atom-character} 逐类计算总位移表示的特征标：

\begin{itemize}
\item \(E\)：5 个原子全部不动，每个三维向量迹为 3，所以 \(\chi_{3N}=15\)；
\item \(C_3\)：转轴穿过 C 和一个 H，这两个原子不动，但 \(1+2\cos(2\pi/3)=0\)，所以 \(\chi_{3N}=0\)；
\item \(C_2\)：只有中心 C 不动，三维 \(\pi\) 旋转迹为 \(1+2\cos\pi=-1\)，所以 \(\chi_{3N}=-1\)；
\item \(S_4\)：只有中心 C 不动，\(S_4\) 在三维矢量表示中的迹为 \(-1\)，所以 \(\chi_{3N}=-1\)；
\item \(\sigma_d\)：镜面包含 C 和两个 H，共 3 个不动原子，每个反射矩阵迹为 1，所以 \(\chi_{3N}=3\)。
\end{itemize}

因此

\begin{equation}
\chi_{3N}
=(15,0,-1,-1,3).
\label{eq:methane-3n-character}
\end{equation}

\(T_d\) 中与本题有关的不可约特征标为

\begin{center}
\begin{tabular}{c|rrrrr}
\toprule
 & \(E\) & \(8C_3\) & \(3C_2\) & \(6S_4\) & \(6\sigma_d\)\\
\midrule
\(A_1\) & 1 & 1 & 1 & 1 & 1\\
\(E\)   & 2 & -1 & 2 & 0 & 0\\
\(T_1\) & 3 & 0 & -1 & 1 & -1\\
\(T_2\) & 3 & 0 & -1 & -1 & 1\\
\bottomrule
\end{tabular}
\end{center}

把 \eqnrefs{eq:methane-3n-character} 代入 \eqnrefs{eq:character-reduction-formula}。例如 \(A_1\) 的重数为

\[
n_{A_1}
=\frac1{24}
[1\cdot15+8\cdot1\cdot0+3\cdot1\cdot(-1)
+6\cdot1\cdot(-1)+6\cdot1\cdot3]
=1.
\]

同理

\[
n_E=1,
\qquad
n_{T_1}=1,
\qquad
n_{T_2}=3,
\]

所以

\begin{equation}
\Gamma_{3N}
=A_1\oplus E\oplus T_1\oplus3T_2.
\label{eq:methane-3n-decomposition}
\end{equation}

三维平移矢量 \((x,y,z)\) 按 \(T_2\) 变换，轴矢量转动 \((R_x,R_y,R_z)\) 按 \(T_1\) 变换，于是

\begin{equation}
\Gamma_{\rm vib}
=A_1\oplus E\oplus2T_2.
\label{eq:methane-vibrational-decomposition}
\end{equation}

维数检查给出 \(1+2+2\times3=9=3N-6\)。这里的“9 个振动自由度”和“4 组基本频率”并不矛盾：\(E\) 是二重简并，两个 \(T_2\) 各自三重简并。

\section{光谱活性也可以从表示直接判断}

选择定则首先是一个一般的群不变量问题。设初态、末态分别按表示 \(\Gamma_i,\Gamma_f\) 变换，跃迁算符 \(O\) 的各分量按 \(\Gamma_O\) 变换。矩阵元
\[
\langle f|O|i\rangle
\]
只有在张量积
\begin{equation}
\Gamma_f^*\otimes\Gamma_O\otimes\Gamma_i
\supset A_1
\label{eq:molecular-general-selection-rule}
\end{equation}
包含全对称表示时才可能非零。等价地，\(\Gamma_f\) 必须出现在 \(\Gamma_O\otimes\Gamma_i\) 的不可约分解中。这一判据同时覆盖红外、Raman、电子跃迁和更高多极跃迁；具体“某一模式是否活性”只是把对应算符表示代入后的特例。

\begin{proof}[解]
设初态、末态、算符分别按表示 $\Gamma_i,\Gamma_f,\Gamma_O$ 变换，基矢记为 $|i,b\rangle$、$|f,a\rangle$，算符独立分量为 $O_\lambda$。全部矩阵元
\begin{equation*}
M_{a\lambda b}=\langle f,a|O_\lambda|i,b\rangle
\end{equation*}
作为指标 $(a,\lambda,b)$ 的集合构成张量
\begin{equation*}
M\in V_f^*\otimes V_O\otimes V_i.
\end{equation*}
若完整哈密顿量在群 $G$ 下不变，跃迁振幅必须在所有对称操作下保持不变。对 $g\in G$，态与算符的变换为
\begin{equation*}
|i,b\rangle\to\sum_{b'}D_i(g)_{b'b}|i,b'\rangle,
\quad
|f,a\rangle\to\sum_{a'}D_f(g)_{a'a}|f,a'\rangle,
\quad
O_\lambda\to\sum_{\lambda'}D_O(g)_{\lambda'\lambda}O_{\lambda'}.
\end{equation*}
由 $\langle f,a|O_\lambda|i,b\rangle$ 的变换规律，$M$ 在群作用下按 $D_f(g)^*\otimes D_O(g)\otimes D_i(g)$ 变换。物理可观测的 $M$ 必须是群不变量，因此对群平均
\begin{equation*}
M = P_{\rm inv}\,M,
\qquad
P_{\rm inv}
= \frac1{|G|}\sum_{g\in G}
 D_f(g)^*\otimes D_O(g)\otimes D_i(g).
\end{equation*}
验证 $P_{\rm inv}$ 为投影：由群乘法性质 $D(gh)=D(g)D(h)$，于是
\begin{equation*}
P_{\rm inv}^2
= \frac1{|G|^2}\sum_{g,h\in G}
 \bigl[D_f(gh)^*\otimes D_O(gh)\otimes D_i(gh)\bigr].
\end{equation*}
固定 $g$ 并对 $h$ 求和，令 $g'=gh$，则 $h$ 遍历 $G$ 时 $g'$ 亦遍历 $G$，因此
\begin{equation*}
P_{\rm inv}^2
= \frac1{|G|^2}\sum_{g\in G}\sum_{g'\in G}
 D_f(g')^*\otimes D_O(g')\otimes D_i(g')
= \frac1{|G|}\sum_{g'\in G}D_f(g')^*\otimes D_O(g')\otimes D_i(g')
= P_{\rm inv},
\end{equation*}
其中用了 $\sum_g=|G|$。因此 $P_{\rm inv}$ 正是张量积表示到全对称子空间 $A_1$ 的投影。由表示论，$P_{\rm inv}\ne0$ 当且仅当
\begin{equation*}
\Gamma_f^*\otimes\Gamma_O\otimes\Gamma_i\supset A_1.
\end{equation*}
若不含 $A_1$，则 $P_{\rm inv}=0$，从而 $M=P_{\rm inv}M=0$，所有相应矩阵元恒为零；若含 $A_1$，群论只说明跃迁允许非零，实际强度仍由约化矩阵元和动力学决定。这便得到\eqnrefs{eq:molecular-general-selection-rule}。
\end{proof}

对振动基态通常有 \(\Gamma_i=A_1\)。电偶极红外跃迁要求振动正常坐标 \(Q_\nu\) 能够使偶极矩的一阶导数

\[
\left(\frac{\partial\boldsymbol\mu}{\partial Q_\nu}\right)_0
\]

非零。因此模式的不可约表示必须出现在矢量表示 \((x,y,z)\) 中。对 \(T_d\)，矢量表示是 \(T_2\)，所以 CH\(_4\) 的两组 \(T_2\) 正规振动为 IR 主动，而 \(A_1,E\) 不由纯电偶极红外一阶激发。

Raman 强度则由极化率张量对正常坐标的一阶导数决定。对称二阶笛卡尔张量的基函数属于

\[
(x^2+y^2+z^2),
\quad
(2z^2-x^2-y^2,x^2-y^2),
\quad
(xy,yz,zx),
\]

在 \(T_d\) 中分别按 \(A_1\oplus E\oplus T_2\) 变换，因此 CH\(_4\) 的 \(A_1,E,T_2\) 模式都可以 Raman 主动。这个选择定则不需要先求出每个正常模式的具体笛卡尔分量。


\section{一般多原子分子的质量加权正常模}

Born--Oppenheimer 势能面 $E_n(\mathbf R)$ 建立以后，振动理论首先是一个一般的多自由度二次型问题，而双原子分子只是只有一个内部振动坐标的特例。设平衡构型为 $\mathbf R^{(0)}$，笛卡尔位移记为 $\delta R_{A\alpha}$。在平衡点附近
\[
E_n(\mathbf R)=E_n(\mathbf R^{(0)})
+\frac12\sum_{A\alpha,B\beta}
K_{A\alpha,B\beta}\delta R_{A\alpha}\delta R_{B\beta}+O(\delta R^3),
\]
其中 Hessian $K$ 的元素量纲为力常数。定义质量加权坐标 $q_{A\alpha}=\sqrt{M_A}\,\delta R_{A\alpha}$ 与质量加权 Hessian
\begin{equation}
F_{A\alpha,B\beta}
=\frac{K_{A\alpha,B\beta}}{\sqrt{M_AM_B}},
\qquad
\sum_{B\beta}F_{A\alpha,B\beta}L_{B\beta,s}=\omega_s^2L_{A\alpha,s}.
\label{eq:mass-weighted-hessian}
\end{equation}
去除整体平移与转动零模以后，剩余正本征值给出分子正常模频率 $\omega_s$。非线性 $N$ 原子分子有 $3N-6$ 个振动模，线性分子有 $3N-5$ 个。

\begin{proof}[解]
孤立分子势能对整体平移不变。将恒等式
\(V(\{\mathbf R_A+\mathbf a\})=V(\{\mathbf R_A\})\)
先对 \(R_{A\alpha}\) 求导，再对 \(a_\beta\) 求导，并令
\(\mathbf a=0\)，便有
\begin{equation}
 \sum_B K_{A\alpha,B\beta}=0 .
\label{eq:hessian-translation-zero}
\end{equation}
所以三个向量
\(t^{(\beta)}_{A\alpha}=\delta_{\alpha\beta}\)
满足 \(Kt^{(\beta)}=0\)。对旋转恒等式
\(V(\{R\mathbf R_A\})=V(\{\mathbf R_A\})\)
作同样微分，在平衡构型处得到
\begin{equation*}
 \sum_{B\beta}K_{A\alpha,B\beta}
 (\boldsymbol\theta\times\mathbf R_B^{(0)})_\beta=0 .
\end{equation*}
非线性构型给出三个独立转动向量；线性构型绕分子轴的向量恒为零，故只有两个。又因为
\(F=M^{-1/2}KM^{-1/2}\)，若 \(Kt=0\)，则
\(F(M^{1/2}t)=0\)。因此质量加权零模应取
\begin{equation*}
 \widetilde t_{A\alpha}^{(\beta)}
 =\sqrt{M_A}\,\delta_{\alpha\beta},\qquad
 \widetilde r_{A\alpha}^{(\beta)}
 =\sqrt{M_A}\,(\mathbf e_\beta\times\mathbf R_A^{(0)})_\alpha .
\end{equation*}
将这些向量正交归一化为列矩阵 \(Z\)，振动子空间的投影算符便是
\(P_{\rm vib}=I-ZZ^{\mathsf T}\)，所需频率由
\(P_{\rm vib}FP_{\rm vib}\) 的正本征值给出。于是
\begin{equation*}
3N-3-3=3N-6\;(\text{非线性}),\qquad
3N-3-2=3N-5\;(\text{线性}).
\end{equation*}
\end{proof}

零模是连续平移与转动对称性的结果，不应与真正的软振动或过渡态的负曲率混淆。分子点群进一步把振动子空间分解为不可约表示；其特征标满足
\(\chi_{\rm vib}=\chi_{3N}-\chi_{\rm trans}-\chi_{\rm rot}\)，从而可在求解 Hessian 之前预言简并度以及红外、拉曼选择定则。若改用键长、键角等内坐标，Wilson \(GF\) 方法会在坐标选择中直接消去刚体运动，与上述投影后的质量加权本征值问题等价。


\section{环形对称性的另一个例子：一般 \texorpdfstring{\(N\)}{N} 元环与苯的 H\"uckel 模型}

考虑 \(N\) 个等价局域轨道 \(|n\rangle\)，只保留最近邻耦合。把在位能量记为 \(\varepsilon_p\)，跳跃积分记为 \(t_\pi\)：

\begin{equation}
H|n\rangle
=\varepsilon_p|n\rangle
+t_\pi(|n+1\rangle+|n-1\rangle),
\qquad n\;\mathrm{mod}\;N.
\label{eq:ring-huckel-hamiltonian}
\end{equation}

定义离散平移 \(T|n\rangle=|n+1\rangle\)。因为 \([H,T]=0\) 且 \(T^N=1\)，\(T\) 的本征值必须满足 \(z^N=1\)，即

\[
z_k=\ee^{-\ii2\pi k/N},
\qquad k=0,\ldots,N-1.
\]

相应归一化 Fourier 态为

\begin{equation}
|k\rangle
=\frac1{\sqrt N}
\sum_{n=0}^{N-1}
 \ee^{\ii2\pi kn/N}|n\rangle.
\label{eq:ring-fourier-state}
\end{equation}

把 \eqnrefs{eq:ring-fourier-state} 代入 \eqnrefs{eq:ring-huckel-hamiltonian}：

\begin{align*}
H|k\rangle
&=\frac1{\sqrt N}\sum_n \ee^{\ii2\pi kn/N}
[\varepsilon_p|n\rangle+t_\pi|n+1\rangle+t_\pi|n-1\rangle]\\
&=\left[\varepsilon_p+t_\pi \ee^{\ii2\pi k/N}
+t_\pi \ee^{-\ii2\pi k/N}\right]|k\rangle,
\end{align*}

因此

\begin{equation}
E_k=\varepsilon_p+2t_\pi\cos\frac{2\pi k}{N}.
\label{eq:ring-huckel-spectrum}
\end{equation}

苯取 \(N=6\)。若 \(t_\pi<0\)，\(k=0\) 为最低成键轨道，\(k=3\) 为最高反键轨道，\(k=\pm1\) 与 \(k=\pm2\) 分别形成两组二重简并。这个例子展示了表示论最实用的一面：先用对称性找到同时对角化群作用的基底，哈密顿量的对角化往往随之变得近乎自动。
